\chapter{Parameter-Efficient Fine-Tuning with LoRA}
\chaptermark{LoRA}
\label{ch:lora}

\section{The problem with full fine-tuning}

Full fine-tuning updates all 809 M parameters of turbo. AdamW stores two extra numbers per parameter (first and second moments),
plus gradients, so memory is several times the model size. It also lets the model drift far from what it learned on 680{,}000 hours,
which risks \emph{catastrophic forgetting} and over-fitting a few hundred hours.

\section{Low-rank adaptation}

\begin{definition}{LoRA (Hu et al., 2021)}{lora}
Take a frozen linear layer $\bm u\mapsto W_0\bm u$ with $W_0\in\R^{d_{\text{out}}\times d_{\text{in}}}$. LoRA replaces it by
\[
\bm u\ \mapsto\ W_0\bm u+\frac{\alpha}{r}\,B A\bm u,\qquad A\in\R^{r\times d_{\text{in}}},\ B\in\R^{d_{\text{out}}\times r},\ r\ll\min(d_{\text{in}},d_{\text{out}}).
\]
Only $A$ and $B$ are trained. Initialisation: $A$ random (Kaiming-uniform), $B=0$. Dropout (here 0.05) is applied to $\bm u$ on the
LoRA branch only.
\end{definition}

\begin{proposition}{Properties of LoRA}{loraprops}
\begin{enumerate}
\item At initialisation the model is exactly the pretrained model.
\item The update $\Delta W=\frac{\alpha}{r}BA$ has $\rank\Delta W\le r$.
\item It needs $r(d_{\text{in}}+d_{\text{out}})$ trainable numbers instead of $d_{\text{in}}d_{\text{out}}$.
\item After training it can be merged, $W=W_0+\frac\alpha rBA$, so inference costs exactly the same as the original model.
\end{enumerate}
\end{proposition}
\begin{proof}
(1) $B=0\Rightarrow\Delta W=0$. (2) $\mathrm{range}(BA)\subseteq\mathrm{range}(B)$, which has dimension at most $r$.
(3) Count entries of $A$ and $B$. (4) $W_0\bm u+\frac\alpha rBA\bm u=(W_0+\frac\alpha rBA)\bm u$ by linearity.
\end{proof}

\begin{proposition}{Gradients and the ``B moves first'' effect}{loragrad}
Let $\bm g=\partial\mathcal L/\partial\bm o$ be the gradient at the layer output $\bm o=W_0\bm u+\frac\alpha rBA\bm u$. Then
\[
\frac{\partial\mathcal L}{\partial B}=\frac\alpha r\,\bm g\,(A\bm u)^\top,\qquad
\frac{\partial\mathcal L}{\partial A}=\frac\alpha r\,B^\top\bm g\,\bm u^\top .
\]
At initialisation ($B=0$) the gradient with respect to $A$ is zero, and only $B$ changes in the first step.
\end{proposition}
\begin{proof}
$\bm o$ depends on $B$ through $\frac\alpha r B\bm v$ with $\bm v=A\bm u$; $\partial(\bm g^\top B\bm v)/\partial B=\bm g\bm v^\top$.
Similarly $\bm g^\top BA\bm u=(B^\top\bm g)^\top A\bm u$ gives $\partial/\partial A=B^\top\bm g\,\bm u^\top$. With $B=0$ this vanishes.
\end{proof}

\begin{intuition}[: why low rank is enough]
Adapting a pretrained model to a related task seems to need only a few new ``directions'' in each layer: the pretrained features are
already rich. The ratio $\alpha/r$ fixes the scale of the update so that changing $r$ does not require re-tuning the learning rate.
Here $\alpha/r=64/32=2$.
\end{intuition}

\section{Where LoRA goes, and the exact parameter count}

LoRA is applied to every attention projection (\texttt{q\_proj, k\_proj, v\_proj, out\_proj}) and both MLP matrices (\texttt{fc1,
fc2}) in every encoder and decoder layer, with $r=32$, $\alpha=64$. In addition the bottom third of encoder layers is frozen:
$\lfloor 32\times0.33\rfloor=10$ layers, so their LoRA matrices are not trained either (lower encoder layers are largely
language-independent acoustic feature extractors; freezing them reduces over-fitting).

\begin{example}{Counting the 86{,}880{,}008 ``trainable'' parameters of Stage 2}{count}
Turbo: $d=1280$, MLP width $5120$. One LoRA pair on a $d_{\text{out}}\times d_{\text{in}}$ matrix costs $32(d_{\text{in}}+d_{\text{out}})$.
\begin{itemize}
\item Attention projection ($1280\to1280$): $32\cdot2560=81{,}920$.\quad \texttt{fc1} or \texttt{fc2}: $32\cdot(1280+5120)=204{,}800$.
\item Encoder layer: $4\cdot81{,}920+2\cdot204{,}800=737{,}280$. Trainable encoder layers: $32-10=22$, giving $16{,}220{,}160$.
\item Decoder layer: self-attention and cross-attention ($8$ projections) $+$ MLP: $8\cdot81{,}920+2\cdot204{,}800=1{,}064{,}960$.
Four layers: $4{,}259{,}840$.
\item LoRA total: $\mathbf{20{,}480{,}000}$.
\item Embedding matrix with the Odia row: $51{,}867\times1280=66{,}389{,}760$ (counted because \texttt{requires\_grad=True}, although a
gradient mask lets only one row of 1{,}280 numbers change; Chapter~\ref{ch:odia}).
\item LID-CTC head: $1280\times8+8=10{,}248$.
\end{itemize}
Sum: $20{,}480{,}000+66{,}389{,}760+10{,}248=86{,}880{,}008$, exactly the 86.9 M printed in the training log. The parameters that
\emph{actually} change are $20{,}480{,}000+1{,}280+10{,}248\approx$ \textbf{20.5 M}, about 2.5\% of the model.
\end{example}

\begin{exercise}{Whisper-small's 13.0 M}{countsmall}
Whisper-small: $d=768$, MLP $3072$, 12 encoder and 12 decoder layers, nothing frozen. Verify the 13.0 M trainable parameters of E1.\par
\textit{Solution.} Attention pair: $32\cdot1536=49{,}152$; MLP pair: $32\cdot3840=122{,}880$. Encoder layer:
$4\cdot49{,}152+2\cdot122{,}880=442{,}368$; decoder layer: $8\cdot49{,}152+2\cdot122{,}880=638{,}976$.
Total $12(442{,}368+638{,}976)=12{,}976{,}128\approx13.0$ M.
\end{exercise}

\section{QLoRA}

QLoRA stores the frozen $W_0$ in 4 bits (the NF4 data type: 16 quantisation levels placed at quantiles of a normal distribution,
which suits roughly Gaussian weights, with per-block scales that are themselves quantised: ``double quantisation''). Matrix products
de-quantise to bfloat16 on the fly. Memory for $W_0$ drops about 4$\times$; LoRA matrices stay in higher precision.

\section{Experiment E1: LoRA vs QLoRA vs full fine-tuning}

Whisper-small, Hindi only, 1{,}000 steps, batch 32; dev CER (FLEURS / IndicVoices / Kathbath). Untrained: 0.49 / 1.39 / 0.72.

\begin{center}
\begin{tabular}{lccccc}
\toprule
Method & trainable & time (min) & FLEURS & IndicVoices & Kathbath\\
\midrule
LoRA ($r=32$) & 13.0 M & 32 & 0.095 & \textbf{0.108} & \textbf{0.125}\\
QLoRA (NF4) & 13.0 M & 76 & \textbf{0.092} & 0.114 & 0.127\\
Full fine-tune (lr $10^{-5}$) & 241.7 M & 28 & 0.101 & 0.173 & 0.143\\
\bottomrule
\end{tabular}
\end{center}

\begin{itemize}
\item QLoRA matches LoRA but is 2.3$\times$ slower (de-quantisation overhead), and memory was not the constraint on a 122 GB machine.
\item Full fine-tuning is clearly worse on the spontaneous IndicVoices set (0.173 vs 0.108), consistent with over-fitting the
read-speech-heavy data and drifting from the pretrained acoustic model. (Caveat: its learning rate was not tuned.)
\end{itemize}
Decision: \textbf{LoRA}.

\begin{incode}
\file{train_whisper.py}: \texttt{LoraConfig(r=32, lora\_alpha=64, lora\_dropout=0.05, target\_modules=[q\_proj, k\_proj, v\_proj,
out\_proj, fc1, fc2])}, then \texttt{--freeze\_enc\_frac 0.33} sets \texttt{requires\_grad=False} on all parameters (including
LoRA) of encoder layers 0--9. \file{load_ft.py} loads the adapter and calls \texttt{merge\_and\_unload()} (Proposition~\ref{prop:loraprops}(4)).
\end{incode}

\begin{vivatip}
Be ready to say why the reported ``87 M trainable'' overstates what is learned: 66.4 M of it is the embedding matrix, of which only
the Odia row can change. The honest figure is about 20.5 M.
\end{vivatip}
